The conceptual consequence centres on the articulation of \(K\),
\(\alpha\), and the electronic section. The reversible reconstruction of
\(K\) preserves data that participate both in positional normalization
and in boundary incidence. The electronic composition uses the
previously constructed coordinates and preserves its centre, transports,
and dimensional realization. The exceptional extension preserves other
relations of the same state under successive representations. The connections
between these objects therefore form part of the result together with
the values of their coordinates.

This organization also determines the order of testing. When
a derivation fixes a value without receiving it as an argument, measurement
tests the consequence obtained. Correspondence with a
physical observable additionally requires identification of the quantity,
normalization, and dimensional realization employed. In this precise sense,
the thesis examines the transition from a parametrized description to an explanation of
the relations determining its parameters. The scope of each result
is that of the proofs presented: structure, equation, and value
constitute inseparable aspects of the determination being tested.

Uniqueness of the record retains two prospective counts and their
respective operations (Section~\ref{act21:sec:relacion-censos}).
The sequence \(40320\to21902\to19446\to79\to1\) culminates in the
heterotypic survivor on the exceptional face; the sequence
\(196\cdot197\cdot196\to331\to2\to1\) culminates in oriented
selection from regional catalogues. The intermediate populations
retain the meaning of each filter. Composition of the selected outputs
with integral inversion connects prospective determination
and reversible recovery.
